\chapter{Curved Spacetime and the Standard Model}

\chapterquote{A field-theory formula is useful only after the smaller calculation inside it is visible.}

This chapter is part of \emph{Non-Abelian Theory, Broken Symmetry, and Advanced Topics}.  The working vocabulary is curved backgrounds, Standard Model, electroweak theory.  The first pass through the chapter should be slow: identify the object being changed, the condition held fixed, and the reason a boundary term or normalization is allowed.  The first concrete theme is reading the gauge group; later sections reuse the same habit in a more compact notation.

For Curved Spacetime and the Standard Model, the calculus-level starting point is tied to reading the gauge group.  Matrices, waves, operators, fields, and gauge variables are introduced only as the calculation requires them.  A formula is accepted after it has been checked in a finite model, differentiated or varied explicitly, and connected to the field-theoretic use that motivates it.

\section{The concrete problem: reading the gauge group}

% QFTFIGURE-BEGIN ch056curvedspacetimeandthestandardmodel-01-gauge-group
\qftfigure{figures/instructional/ch056curvedspacetimeandthestandardmodel/ch056curvedspacetimeandthestandardmodel-01-gauge-group.pdf}{The three linked sectors show the gauge-group structure as connected but distinct pieces of one theory.}{fig:ch056curvedspacetimeandthestandardmodel-01-gauge-group}
% QFTFIGURE-END ch056curvedspacetimeandthestandardmodel-01-gauge-group












The work of this section is reading the gauge group.  In the chapter heading the vocabulary is curved backgrounds, Standard Model, electroweak theory, but the first objects on the page are more prosaic: gauge groups, matter multiplets, Yukawa couplings, electroweak breaking, curved backgrounds, and effective descriptions.  The formal notation will be useful only after it has been tied to a limit, a symmetry, a normalization condition, or an integration-by-parts step.  In Curved Spacetime and the Standard Model, section 1, the useful habit is to name the object, the operation, and the assumption before using the compact notation.

The finite model is a table listing fields, representations, and charges, checked term by term for gauge invariance.  In that model no mystery is available: every derivative is an ordinary derivative with respect to one listed variable, every inner product is a sum, and every boundary assumption can be written at the endpoints.  This is why the finite model is not a toy in the dismissive sense.  It is the bookkeeping device that prevents continuum notation from becoming a fog of indices.  For Curved Spacetime and the Standard Model, section 1, that bookkeeping is deliberately modest, but it is what later keeps signs, factors of \(2\pi\), and normalization constants from turning into guesses.

The central idea for this chapter is the Standard Model Lagrangian is the collection of lowest-dimension local terms allowed by its gauge symmetries and field content.  To test that idea, strip the formula down until only the operation remains.  A sign change after raising or lowering an index means the metric has entered.  A derivative moving from one factor to another means a boundary term has been handled.  A normalization constant means that some unit object, such as a delta function, basis vector, or one-particle state, has been fixed.  Read in the setting of Curved Spacetime and the Standard Model, section 1, the line is a check on the calculation rather than an invitation to memorize another rule.

For the concrete problem: reading the gauge group, the calculation should also pass a dimensional check.  The left and right sides must carry the same units; more subtly, they must transform in the same way under the symmetries already introduced.  This is not a proof, but it is a quick way to catch wrong factors of \(2\pi\), signs in the metric, missing powers of the lattice spacing, or an omitted charge.  The same step will look more abstract later; in Curved Spacetime and the Standard Model, section 1 it is kept close to the finite calculation so the limiting move remains visible.

The bridge to quantum field theory is direct.  The final application uses every earlier tool: spinors, gauge fields, symmetry breaking, renormalization, and effective reasoning.  Later notation will carry more indices and fewer verbal reminders, so the safe habit is to keep asking which operation is being performed and which quantity is being held fixed.  At this point in Curved Spacetime and the Standard Model, section 1, it is worth asking what would fail if the boundary condition, normalization, or symmetry assumption were changed.

One warning is worth making explicit here.  A compact Lagrangian hides many definitions; each symbol carries representation, charge, and normalization data.  This warning points to the delicate step of the derivation, not to a decorative aside.  In Curved Spacetime and the Standard Model, section 1, the calculation is meant to leave a trail from an ordinary derivative, sum, matrix product, or Gaussian integral to the field-theory formula.

The section closes by keeping reading the gauge group connected to Curved Spacetime and the Standard Model.  The same general operation will be used with different objects in neighboring chapters, so the present calculation is both a result and a rehearsal.  In Curved Spacetime and the Standard Model, section 1, the useful habit is to name the object, the operation, and the assumption before using the compact notation.



\paragraph{Step-by-step derivation.}

For Curved Spacetime and the Standard Model: The concrete problem: reading the gauge group, the derivation is written from the smallest usable model upward.

List the fields with their representations under \(SU(3)_c\times SU(2)_L\times U(1)_Y\).  A term in the Lagrangian is allowed only if all gauge indices can be contracted to a singlet and the hypercharges sum to zero.  For example, the charged-lepton Yukawa term has the schematic form
\[
  -y_e\bar L\Phi e_R+\hbox{h.c.}
\]
The \(SU(2)\) indices of \(L\) and \(\Phi\) are contracted, and the hypercharges are chosen so the product is neutral.  When \(\Phi\) is expanded around its vacuum value, the same term becomes a fermion mass term.  This single check uses representations, gauge invariance, spontaneous symmetry breaking, and ordinary dimensional analysis.  In Curved Spacetime and the Standard Model, section 1, the useful habit is to name the object, the operation, and the assumption before using the compact notation.

The formulas to keep in view are

\[
  G_{\mathrm{SM}}=SU(3)_c\times SU(2)_L\times U(1)_Y
\]
\[
  D_\mu=\p_\mu+\ii g_sG_\mu^aT^a+\ii gW_\mu^i\tau^i+\ii g'YB_\mu
\]
\[
  \Lag\supset -y_e\bar L\Phi e_R+\hbox{h.c.}
\]

In this setting the calculation for Curved Spacetime and the Standard Model: The concrete problem: reading the gauge group is complete when the finite model, the limiting step, and the normalization or boundary condition can all be named without looking back at the prose.




\begin{example}[A worked calculation for Curved Spacetime and the Standard Model: The concrete problem: reading the gauge group]
Check a charged-lepton Yukawa term.  The left-handed lepton \(L\) is an \(SU(2)\) doublet, the Higgs field \(\Phi\) is a doublet, and \(e_R\) is a singlet.  Contract the two doublet indices and choose hypercharges so the sum is zero.  After \(\Phi\) gets a vacuum value, the term becomes \(m_e\bar e e\) with \(m_e=y_ev/\sqrt2\).  In Curved Spacetime and the Standard Model, section 1, the useful habit is to name the object, the operation, and the assumption before using the compact notation.
\end{example}



\begin{exercise}
State the Standard Model gauge group.  Use the notation of Curved Spacetime and the Standard Model: The concrete problem: reading the gauge group.
\begin{answer}
\(SU(3)_c\times SU(2)_L\times U(1)_Y\).
\end{answer}
\end{exercise}

\begin{exercise}
Why must hypercharges in an interaction term sum to zero?  Use the notation of Curved Spacetime and the Standard Model: The concrete problem: reading the gauge group.
\begin{answer}
Otherwise the term picks up a local \(U(1)_Y\) phase and is not gauge invariant.
\end{answer}
\end{exercise}



\section{Objects, notation, and units: assigning representations}

% QFTFIGURE-BEGIN ch056curvedspacetimeandthestandardmodel-02-representations
\qftfigure{figures/instructional/ch056curvedspacetimeandthestandardmodel/ch056curvedspacetimeandthestandardmodel-02-representations.pdf}{The grouped multiplets visualize matter fields as organized by their symmetry charges.}{fig:ch056curvedspacetimeandthestandardmodel-02-representations}
% QFTFIGURE-END ch056curvedspacetimeandthestandardmodel-02-representations












The work of this section is assigning representations.  In the chapter heading the vocabulary is curved backgrounds, Standard Model, electroweak theory, but the first objects on the page are more prosaic: gauge groups, matter multiplets, Yukawa couplings, electroweak breaking, curved backgrounds, and effective descriptions.  The formal notation will be useful only after it has been tied to a limit, a symmetry, a normalization condition, or an integration-by-parts step.  In Curved Spacetime and the Standard Model, section 2, the useful habit is to name the object, the operation, and the assumption before using the compact notation.

The finite model is a table listing fields, representations, and charges, checked term by term for gauge invariance.  In that model no mystery is available: every derivative is an ordinary derivative with respect to one listed variable, every inner product is a sum, and every boundary assumption can be written at the endpoints.  This is why the finite model is not a toy in the dismissive sense.  It is the bookkeeping device that prevents continuum notation from becoming a fog of indices.  For Curved Spacetime and the Standard Model, section 2, that bookkeeping is deliberately modest, but it is what later keeps signs, factors of \(2\pi\), and normalization constants from turning into guesses.

The central idea for this chapter is the Standard Model Lagrangian is the collection of lowest-dimension local terms allowed by its gauge symmetries and field content.  To test that idea, strip the formula down until only the operation remains.  A sign change after raising or lowering an index means the metric has entered.  A derivative moving from one factor to another means a boundary term has been handled.  A normalization constant means that some unit object, such as a delta function, basis vector, or one-particle state, has been fixed.  Read in the setting of Curved Spacetime and the Standard Model, section 2, the line is a check on the calculation rather than an invitation to memorize another rule.

For objects, notation, and units: assigning representations, the calculation should also pass a dimensional check.  The left and right sides must carry the same units; more subtly, they must transform in the same way under the symmetries already introduced.  This is not a proof, but it is a quick way to catch wrong factors of \(2\pi\), signs in the metric, missing powers of the lattice spacing, or an omitted charge.  The same step will look more abstract later; in Curved Spacetime and the Standard Model, section 2 it is kept close to the finite calculation so the limiting move remains visible.

The bridge to quantum field theory is direct.  The final application uses every earlier tool: spinors, gauge fields, symmetry breaking, renormalization, and effective reasoning.  Later notation will carry more indices and fewer verbal reminders, so the safe habit is to keep asking which operation is being performed and which quantity is being held fixed.  At this point in Curved Spacetime and the Standard Model, section 2, it is worth asking what would fail if the boundary condition, normalization, or symmetry assumption were changed.

One warning is worth making explicit here.  A compact Lagrangian hides many definitions; each symbol carries representation, charge, and normalization data.  This warning points to the delicate step of the derivation, not to a decorative aside.  In Curved Spacetime and the Standard Model, section 2, the calculation is meant to leave a trail from an ordinary derivative, sum, matrix product, or Gaussian integral to the field-theory formula.

The section closes by keeping assigning representations connected to Curved Spacetime and the Standard Model.  The same general operation will be used with different objects in neighboring chapters, so the present calculation is both a result and a rehearsal.  In Curved Spacetime and the Standard Model, section 2, the useful habit is to name the object, the operation, and the assumption before using the compact notation.



\paragraph{Step-by-step derivation.}

For Curved Spacetime and the Standard Model: Objects, notation, and units: assigning representations, the derivation is written from the smallest usable model upward.

List the fields with their representations under \(SU(3)_c\times SU(2)_L\times U(1)_Y\).  A term in the Lagrangian is allowed only if all gauge indices can be contracted to a singlet and the hypercharges sum to zero.  For example, the charged-lepton Yukawa term has the schematic form
\[
  -y_e\bar L\Phi e_R+\hbox{h.c.}
\]
The \(SU(2)\) indices of \(L\) and \(\Phi\) are contracted, and the hypercharges are chosen so the product is neutral.  When \(\Phi\) is expanded around its vacuum value, the same term becomes a fermion mass term.  This single check uses representations, gauge invariance, spontaneous symmetry breaking, and ordinary dimensional analysis.  In Curved Spacetime and the Standard Model, section 2, the useful habit is to name the object, the operation, and the assumption before using the compact notation.

The formulas to keep in view are

\[
  G_{\mathrm{SM}}=SU(3)_c\times SU(2)_L\times U(1)_Y
\]
\[
  D_\mu=\p_\mu+\ii g_sG_\mu^aT^a+\ii gW_\mu^i\tau^i+\ii g'YB_\mu
\]
\[
  \Lag\supset -y_e\bar L\Phi e_R+\hbox{h.c.}
\]

In this setting the calculation for Curved Spacetime and the Standard Model: Objects, notation, and units: assigning representations is complete when the finite model, the limiting step, and the normalization or boundary condition can all be named without looking back at the prose.




\begin{example}[A worked calculation for Curved Spacetime and the Standard Model: Objects, notation, and units: assigning representations]
Check a charged-lepton Yukawa term.  The left-handed lepton \(L\) is an \(SU(2)\) doublet, the Higgs field \(\Phi\) is a doublet, and \(e_R\) is a singlet.  Contract the two doublet indices and choose hypercharges so the sum is zero.  After \(\Phi\) gets a vacuum value, the term becomes \(m_e\bar e e\) with \(m_e=y_ev/\sqrt2\).  In Curved Spacetime and the Standard Model, section 2, the useful habit is to name the object, the operation, and the assumption before using the compact notation.
\end{example}



\begin{exercise}
State the Standard Model gauge group.  Use the notation of Curved Spacetime and the Standard Model: Objects, notation, and units: assigning representations.
\begin{answer}
\(SU(3)_c\times SU(2)_L\times U(1)_Y\).
\end{answer}
\end{exercise}

\begin{exercise}
Why must hypercharges in an interaction term sum to zero?  Use the notation of Curved Spacetime and the Standard Model: Objects, notation, and units: assigning representations.
\begin{answer}
Otherwise the term picks up a local \(U(1)_Y\) phase and is not gauge invariant.
\end{answer}
\end{exercise}



\section{The finite calculation: building covariant derivatives}

% QFTFIGURE-BEGIN ch056curvedspacetimeandthestandardmodel-03-yukawa
\qftfigure{figures/instructional/ch056curvedspacetimeandthestandardmodel/ch056curvedspacetimeandthestandardmodel-03-yukawa.pdf}{The fermion pair meeting a scalar line shows how masses arise after symmetry breaking.}{fig:ch056curvedspacetimeandthestandardmodel-03-yukawa}
% QFTFIGURE-END ch056curvedspacetimeandthestandardmodel-03-yukawa












The work of this section is building covariant derivatives.  In the chapter heading the vocabulary is curved backgrounds, Standard Model, electroweak theory, but the first objects on the page are more prosaic: gauge groups, matter multiplets, Yukawa couplings, electroweak breaking, curved backgrounds, and effective descriptions.  The formal notation will be useful only after it has been tied to a limit, a symmetry, a normalization condition, or an integration-by-parts step.  In Curved Spacetime and the Standard Model, section 3, the useful habit is to name the object, the operation, and the assumption before using the compact notation.

The finite model is a table listing fields, representations, and charges, checked term by term for gauge invariance.  In that model no mystery is available: every derivative is an ordinary derivative with respect to one listed variable, every inner product is a sum, and every boundary assumption can be written at the endpoints.  This is why the finite model is not a toy in the dismissive sense.  It is the bookkeeping device that prevents continuum notation from becoming a fog of indices.  For Curved Spacetime and the Standard Model, section 3, that bookkeeping is deliberately modest, but it is what later keeps signs, factors of \(2\pi\), and normalization constants from turning into guesses.

The central idea for this chapter is the Standard Model Lagrangian is the collection of lowest-dimension local terms allowed by its gauge symmetries and field content.  To test that idea, strip the formula down until only the operation remains.  A sign change after raising or lowering an index means the metric has entered.  A derivative moving from one factor to another means a boundary term has been handled.  A normalization constant means that some unit object, such as a delta function, basis vector, or one-particle state, has been fixed.  Read in the setting of Curved Spacetime and the Standard Model, section 3, the line is a check on the calculation rather than an invitation to memorize another rule.

For the finite calculation: building covariant derivatives, the calculation should also pass a dimensional check.  The left and right sides must carry the same units; more subtly, they must transform in the same way under the symmetries already introduced.  This is not a proof, but it is a quick way to catch wrong factors of \(2\pi\), signs in the metric, missing powers of the lattice spacing, or an omitted charge.  The same step will look more abstract later; in Curved Spacetime and the Standard Model, section 3 it is kept close to the finite calculation so the limiting move remains visible.

The bridge to quantum field theory is direct.  The final application uses every earlier tool: spinors, gauge fields, symmetry breaking, renormalization, and effective reasoning.  Later notation will carry more indices and fewer verbal reminders, so the safe habit is to keep asking which operation is being performed and which quantity is being held fixed.  At this point in Curved Spacetime and the Standard Model, section 3, it is worth asking what would fail if the boundary condition, normalization, or symmetry assumption were changed.

One warning is worth making explicit here.  A compact Lagrangian hides many definitions; each symbol carries representation, charge, and normalization data.  This warning points to the delicate step of the derivation, not to a decorative aside.  In Curved Spacetime and the Standard Model, section 3, the calculation is meant to leave a trail from an ordinary derivative, sum, matrix product, or Gaussian integral to the field-theory formula.

The section closes by keeping building covariant derivatives connected to Curved Spacetime and the Standard Model.  The same general operation will be used with different objects in neighboring chapters, so the present calculation is both a result and a rehearsal.  In Curved Spacetime and the Standard Model, section 3, the useful habit is to name the object, the operation, and the assumption before using the compact notation.



\paragraph{Step-by-step derivation.}

For Curved Spacetime and the Standard Model: The finite calculation: building covariant derivatives, the derivation is written from the smallest usable model upward.

List the fields with their representations under \(SU(3)_c\times SU(2)_L\times U(1)_Y\).  A term in the Lagrangian is allowed only if all gauge indices can be contracted to a singlet and the hypercharges sum to zero.  For example, the charged-lepton Yukawa term has the schematic form
\[
  -y_e\bar L\Phi e_R+\hbox{h.c.}
\]
The \(SU(2)\) indices of \(L\) and \(\Phi\) are contracted, and the hypercharges are chosen so the product is neutral.  When \(\Phi\) is expanded around its vacuum value, the same term becomes a fermion mass term.  This single check uses representations, gauge invariance, spontaneous symmetry breaking, and ordinary dimensional analysis.  In Curved Spacetime and the Standard Model, section 3, the useful habit is to name the object, the operation, and the assumption before using the compact notation.

The formulas to keep in view are

\[
  G_{\mathrm{SM}}=SU(3)_c\times SU(2)_L\times U(1)_Y
\]
\[
  D_\mu=\p_\mu+\ii g_sG_\mu^aT^a+\ii gW_\mu^i\tau^i+\ii g'YB_\mu
\]
\[
  \Lag\supset -y_e\bar L\Phi e_R+\hbox{h.c.}
\]

In this setting the calculation for Curved Spacetime and the Standard Model: The finite calculation: building covariant derivatives is complete when the finite model, the limiting step, and the normalization or boundary condition can all be named without looking back at the prose.




\begin{example}[A worked calculation for Curved Spacetime and the Standard Model: The finite calculation: building covariant derivatives]
Check a charged-lepton Yukawa term.  The left-handed lepton \(L\) is an \(SU(2)\) doublet, the Higgs field \(\Phi\) is a doublet, and \(e_R\) is a singlet.  Contract the two doublet indices and choose hypercharges so the sum is zero.  After \(\Phi\) gets a vacuum value, the term becomes \(m_e\bar e e\) with \(m_e=y_ev/\sqrt2\).  In Curved Spacetime and the Standard Model, section 3, the useful habit is to name the object, the operation, and the assumption before using the compact notation.
\end{example}



\begin{exercise}
State the Standard Model gauge group.  Use the notation of Curved Spacetime and the Standard Model: The finite calculation: building covariant derivatives.
\begin{answer}
\(SU(3)_c\times SU(2)_L\times U(1)_Y\).
\end{answer}
\end{exercise}

\begin{exercise}
Why must hypercharges in an interaction term sum to zero?  Use the notation of Curved Spacetime and the Standard Model: The finite calculation: building covariant derivatives.
\begin{answer}
Otherwise the term picks up a local \(U(1)_Y\) phase and is not gauge invariant.
\end{answer}
\end{exercise}



\section{The central derivation: checking Yukawa terms}

% QFTFIGURE-BEGIN ch056curvedspacetimeandthestandardmodel-04-electroweak
\qftfigure{figures/instructional/ch056curvedspacetimeandthestandardmodel/ch056curvedspacetimeandthestandardmodel-04-electroweak.pdf}{The mixing diagram shows weak and hypercharge fields reorganizing into physical fields.}{fig:ch056curvedspacetimeandthestandardmodel-04-electroweak}
% QFTFIGURE-END ch056curvedspacetimeandthestandardmodel-04-electroweak












The work of this section is checking Yukawa terms.  In the chapter heading the vocabulary is curved backgrounds, Standard Model, electroweak theory, but the first objects on the page are more prosaic: gauge groups, matter multiplets, Yukawa couplings, electroweak breaking, curved backgrounds, and effective descriptions.  The formal notation will be useful only after it has been tied to a limit, a symmetry, a normalization condition, or an integration-by-parts step.  In Curved Spacetime and the Standard Model, section 4, the useful habit is to name the object, the operation, and the assumption before using the compact notation.

The finite model is a table listing fields, representations, and charges, checked term by term for gauge invariance.  In that model no mystery is available: every derivative is an ordinary derivative with respect to one listed variable, every inner product is a sum, and every boundary assumption can be written at the endpoints.  This is why the finite model is not a toy in the dismissive sense.  It is the bookkeeping device that prevents continuum notation from becoming a fog of indices.  For Curved Spacetime and the Standard Model, section 4, that bookkeeping is deliberately modest, but it is what later keeps signs, factors of \(2\pi\), and normalization constants from turning into guesses.

The central idea for this chapter is the Standard Model Lagrangian is the collection of lowest-dimension local terms allowed by its gauge symmetries and field content.  To test that idea, strip the formula down until only the operation remains.  A sign change after raising or lowering an index means the metric has entered.  A derivative moving from one factor to another means a boundary term has been handled.  A normalization constant means that some unit object, such as a delta function, basis vector, or one-particle state, has been fixed.  Read in the setting of Curved Spacetime and the Standard Model, section 4, the line is a check on the calculation rather than an invitation to memorize another rule.

For the central derivation: checking yukawa terms, the calculation should also pass a dimensional check.  The left and right sides must carry the same units; more subtly, they must transform in the same way under the symmetries already introduced.  This is not a proof, but it is a quick way to catch wrong factors of \(2\pi\), signs in the metric, missing powers of the lattice spacing, or an omitted charge.  The same step will look more abstract later; in Curved Spacetime and the Standard Model, section 4 it is kept close to the finite calculation so the limiting move remains visible.

The bridge to quantum field theory is direct.  The final application uses every earlier tool: spinors, gauge fields, symmetry breaking, renormalization, and effective reasoning.  Later notation will carry more indices and fewer verbal reminders, so the safe habit is to keep asking which operation is being performed and which quantity is being held fixed.  At this point in Curved Spacetime and the Standard Model, section 4, it is worth asking what would fail if the boundary condition, normalization, or symmetry assumption were changed.

One warning is worth making explicit here.  A compact Lagrangian hides many definitions; each symbol carries representation, charge, and normalization data.  This warning points to the delicate step of the derivation, not to a decorative aside.  In Curved Spacetime and the Standard Model, section 4, the calculation is meant to leave a trail from an ordinary derivative, sum, matrix product, or Gaussian integral to the field-theory formula.

The section closes by keeping checking Yukawa terms connected to Curved Spacetime and the Standard Model.  The same general operation will be used with different objects in neighboring chapters, so the present calculation is both a result and a rehearsal.  In Curved Spacetime and the Standard Model, section 4, the useful habit is to name the object, the operation, and the assumption before using the compact notation.



\paragraph{Step-by-step derivation.}

For Curved Spacetime and the Standard Model: The central derivation: checking Yukawa terms, the derivation is written from the smallest usable model upward.

List the fields with their representations under \(SU(3)_c\times SU(2)_L\times U(1)_Y\).  A term in the Lagrangian is allowed only if all gauge indices can be contracted to a singlet and the hypercharges sum to zero.  For example, the charged-lepton Yukawa term has the schematic form
\[
  -y_e\bar L\Phi e_R+\hbox{h.c.}
\]
The \(SU(2)\) indices of \(L\) and \(\Phi\) are contracted, and the hypercharges are chosen so the product is neutral.  When \(\Phi\) is expanded around its vacuum value, the same term becomes a fermion mass term.  This single check uses representations, gauge invariance, spontaneous symmetry breaking, and ordinary dimensional analysis.  In Curved Spacetime and the Standard Model, section 4, the useful habit is to name the object, the operation, and the assumption before using the compact notation.

The formulas to keep in view are

\[
  G_{\mathrm{SM}}=SU(3)_c\times SU(2)_L\times U(1)_Y
\]
\[
  D_\mu=\p_\mu+\ii g_sG_\mu^aT^a+\ii gW_\mu^i\tau^i+\ii g'YB_\mu
\]
\[
  \Lag\supset -y_e\bar L\Phi e_R+\hbox{h.c.}
\]

In this setting the calculation for Curved Spacetime and the Standard Model: The central derivation: checking Yukawa terms is complete when the finite model, the limiting step, and the normalization or boundary condition can all be named without looking back at the prose.




\begin{example}[A worked calculation for Curved Spacetime and the Standard Model: The central derivation: checking Yukawa terms]
Check a charged-lepton Yukawa term.  The left-handed lepton \(L\) is an \(SU(2)\) doublet, the Higgs field \(\Phi\) is a doublet, and \(e_R\) is a singlet.  Contract the two doublet indices and choose hypercharges so the sum is zero.  After \(\Phi\) gets a vacuum value, the term becomes \(m_e\bar e e\) with \(m_e=y_ev/\sqrt2\).  In Curved Spacetime and the Standard Model, section 4, the useful habit is to name the object, the operation, and the assumption before using the compact notation.
\end{example}



\begin{exercise}
State the Standard Model gauge group.  Use the notation of Curved Spacetime and the Standard Model: The central derivation: checking Yukawa terms.
\begin{answer}
\(SU(3)_c\times SU(2)_L\times U(1)_Y\).
\end{answer}
\end{exercise}

\begin{exercise}
Why must hypercharges in an interaction term sum to zero?  Use the notation of Curved Spacetime and the Standard Model: The central derivation: checking Yukawa terms.
\begin{answer}
Otherwise the term picks up a local \(U(1)_Y\) phase and is not gauge invariant.
\end{answer}
\end{exercise}



\section{Worked calculation: breaking electroweak symmetry}

% QFTFIGURE-BEGIN ch056curvedspacetimeandthestandardmodel-05-curved-background
\qftfigure{figures/instructional/ch056curvedspacetimeandthestandardmodel/ch056curvedspacetimeandthestandardmodel-05-curved-background.pdf}{The field pattern drawn over a curved grid makes clear that the background geometry changes the bookkeeping of local fields.}{fig:ch056curvedspacetimeandthestandardmodel-05-curved-background}
% QFTFIGURE-END ch056curvedspacetimeandthestandardmodel-05-curved-background












The work of this section is breaking electroweak symmetry.  In the chapter heading the vocabulary is curved backgrounds, Standard Model, electroweak theory, but the first objects on the page are more prosaic: gauge groups, matter multiplets, Yukawa couplings, electroweak breaking, curved backgrounds, and effective descriptions.  The formal notation will be useful only after it has been tied to a limit, a symmetry, a normalization condition, or an integration-by-parts step.  In Curved Spacetime and the Standard Model, section 5, the useful habit is to name the object, the operation, and the assumption before using the compact notation.

The finite model is a table listing fields, representations, and charges, checked term by term for gauge invariance.  In that model no mystery is available: every derivative is an ordinary derivative with respect to one listed variable, every inner product is a sum, and every boundary assumption can be written at the endpoints.  This is why the finite model is not a toy in the dismissive sense.  It is the bookkeeping device that prevents continuum notation from becoming a fog of indices.  For Curved Spacetime and the Standard Model, section 5, that bookkeeping is deliberately modest, but it is what later keeps signs, factors of \(2\pi\), and normalization constants from turning into guesses.

The central idea for this chapter is the Standard Model Lagrangian is the collection of lowest-dimension local terms allowed by its gauge symmetries and field content.  To test that idea, strip the formula down until only the operation remains.  A sign change after raising or lowering an index means the metric has entered.  A derivative moving from one factor to another means a boundary term has been handled.  A normalization constant means that some unit object, such as a delta function, basis vector, or one-particle state, has been fixed.  Read in the setting of Curved Spacetime and the Standard Model, section 5, the line is a check on the calculation rather than an invitation to memorize another rule.

For worked calculation: breaking electroweak symmetry, the calculation should also pass a dimensional check.  The left and right sides must carry the same units; more subtly, they must transform in the same way under the symmetries already introduced.  This is not a proof, but it is a quick way to catch wrong factors of \(2\pi\), signs in the metric, missing powers of the lattice spacing, or an omitted charge.  The same step will look more abstract later; in Curved Spacetime and the Standard Model, section 5 it is kept close to the finite calculation so the limiting move remains visible.

The bridge to quantum field theory is direct.  The final application uses every earlier tool: spinors, gauge fields, symmetry breaking, renormalization, and effective reasoning.  Later notation will carry more indices and fewer verbal reminders, so the safe habit is to keep asking which operation is being performed and which quantity is being held fixed.  At this point in Curved Spacetime and the Standard Model, section 5, it is worth asking what would fail if the boundary condition, normalization, or symmetry assumption were changed.

One warning is worth making explicit here.  A compact Lagrangian hides many definitions; each symbol carries representation, charge, and normalization data.  This warning points to the delicate step of the derivation, not to a decorative aside.  In Curved Spacetime and the Standard Model, section 5, the calculation is meant to leave a trail from an ordinary derivative, sum, matrix product, or Gaussian integral to the field-theory formula.

The section closes by keeping breaking electroweak symmetry connected to Curved Spacetime and the Standard Model.  The same general operation will be used with different objects in neighboring chapters, so the present calculation is both a result and a rehearsal.  In Curved Spacetime and the Standard Model, section 5, the useful habit is to name the object, the operation, and the assumption before using the compact notation.



\paragraph{Step-by-step derivation.}

For Curved Spacetime and the Standard Model: Worked calculation: breaking electroweak symmetry, the derivation is written from the smallest usable model upward.

List the fields with their representations under \(SU(3)_c\times SU(2)_L\times U(1)_Y\).  A term in the Lagrangian is allowed only if all gauge indices can be contracted to a singlet and the hypercharges sum to zero.  For example, the charged-lepton Yukawa term has the schematic form
\[
  -y_e\bar L\Phi e_R+\hbox{h.c.}
\]
The \(SU(2)\) indices of \(L\) and \(\Phi\) are contracted, and the hypercharges are chosen so the product is neutral.  When \(\Phi\) is expanded around its vacuum value, the same term becomes a fermion mass term.  This single check uses representations, gauge invariance, spontaneous symmetry breaking, and ordinary dimensional analysis.  In Curved Spacetime and the Standard Model, section 5, the useful habit is to name the object, the operation, and the assumption before using the compact notation.

The formulas to keep in view are

\[
  G_{\mathrm{SM}}=SU(3)_c\times SU(2)_L\times U(1)_Y
\]
\[
  D_\mu=\p_\mu+\ii g_sG_\mu^aT^a+\ii gW_\mu^i\tau^i+\ii g'YB_\mu
\]
\[
  \Lag\supset -y_e\bar L\Phi e_R+\hbox{h.c.}
\]

In this setting the calculation for Curved Spacetime and the Standard Model: Worked calculation: breaking electroweak symmetry is complete when the finite model, the limiting step, and the normalization or boundary condition can all be named without looking back at the prose.




\begin{example}[A worked calculation for Curved Spacetime and the Standard Model: Worked calculation: breaking electroweak symmetry]
Check a charged-lepton Yukawa term.  The left-handed lepton \(L\) is an \(SU(2)\) doublet, the Higgs field \(\Phi\) is a doublet, and \(e_R\) is a singlet.  Contract the two doublet indices and choose hypercharges so the sum is zero.  After \(\Phi\) gets a vacuum value, the term becomes \(m_e\bar e e\) with \(m_e=y_ev/\sqrt2\).  In Curved Spacetime and the Standard Model, section 5, the useful habit is to name the object, the operation, and the assumption before using the compact notation.
\end{example}



\begin{exercise}
State the Standard Model gauge group.  Use the notation of Curved Spacetime and the Standard Model: Worked calculation: breaking electroweak symmetry.
\begin{answer}
\(SU(3)_c\times SU(2)_L\times U(1)_Y\).
\end{answer}
\end{exercise}

\begin{exercise}
Why must hypercharges in an interaction term sum to zero?  Use the notation of Curved Spacetime and the Standard Model: Worked calculation: breaking electroweak symmetry.
\begin{answer}
Otherwise the term picks up a local \(U(1)_Y\) phase and is not gauge invariant.
\end{answer}
\end{exercise}



\section{Boundary conditions, normalization, and signs: identifying physical fields}

% QFTFIGURE-BEGIN ch056curvedspacetimeandthestandardmodel-06-course-map
\qftfigure{figures/instructional/ch056curvedspacetimeandthestandardmodel/ch056curvedspacetimeandthestandardmodel-06-course-map.pdf}{The network map ties calculus, fields, symmetry, renormalization, and the Standard Model into one path through the book.}{fig:ch056curvedspacetimeandthestandardmodel-06-course-map}
% QFTFIGURE-END ch056curvedspacetimeandthestandardmodel-06-course-map












The work of this section is identifying physical fields.  In the chapter heading the vocabulary is curved backgrounds, Standard Model, electroweak theory, but the first objects on the page are more prosaic: gauge groups, matter multiplets, Yukawa couplings, electroweak breaking, curved backgrounds, and effective descriptions.  The formal notation will be useful only after it has been tied to a limit, a symmetry, a normalization condition, or an integration-by-parts step.  In Curved Spacetime and the Standard Model, section 6, the useful habit is to name the object, the operation, and the assumption before using the compact notation.

The finite model is a table listing fields, representations, and charges, checked term by term for gauge invariance.  In that model no mystery is available: every derivative is an ordinary derivative with respect to one listed variable, every inner product is a sum, and every boundary assumption can be written at the endpoints.  This is why the finite model is not a toy in the dismissive sense.  It is the bookkeeping device that prevents continuum notation from becoming a fog of indices.  For Curved Spacetime and the Standard Model, section 6, that bookkeeping is deliberately modest, but it is what later keeps signs, factors of \(2\pi\), and normalization constants from turning into guesses.

The central idea for this chapter is the Standard Model Lagrangian is the collection of lowest-dimension local terms allowed by its gauge symmetries and field content.  To test that idea, strip the formula down until only the operation remains.  A sign change after raising or lowering an index means the metric has entered.  A derivative moving from one factor to another means a boundary term has been handled.  A normalization constant means that some unit object, such as a delta function, basis vector, or one-particle state, has been fixed.  Read in the setting of Curved Spacetime and the Standard Model, section 6, the line is a check on the calculation rather than an invitation to memorize another rule.

For boundary conditions, normalization, and signs: identifying physical fields, the calculation should also pass a dimensional check.  The left and right sides must carry the same units; more subtly, they must transform in the same way under the symmetries already introduced.  This is not a proof, but it is a quick way to catch wrong factors of \(2\pi\), signs in the metric, missing powers of the lattice spacing, or an omitted charge.  The same step will look more abstract later; in Curved Spacetime and the Standard Model, section 6 it is kept close to the finite calculation so the limiting move remains visible.

The bridge to quantum field theory is direct.  The final application uses every earlier tool: spinors, gauge fields, symmetry breaking, renormalization, and effective reasoning.  Later notation will carry more indices and fewer verbal reminders, so the safe habit is to keep asking which operation is being performed and which quantity is being held fixed.  At this point in Curved Spacetime and the Standard Model, section 6, it is worth asking what would fail if the boundary condition, normalization, or symmetry assumption were changed.

One warning is worth making explicit here.  A compact Lagrangian hides many definitions; each symbol carries representation, charge, and normalization data.  This warning points to the delicate step of the derivation, not to a decorative aside.  In Curved Spacetime and the Standard Model, section 6, the calculation is meant to leave a trail from an ordinary derivative, sum, matrix product, or Gaussian integral to the field-theory formula.

The section closes by keeping identifying physical fields connected to Curved Spacetime and the Standard Model.  The same general operation will be used with different objects in neighboring chapters, so the present calculation is both a result and a rehearsal.  In Curved Spacetime and the Standard Model, section 6, the useful habit is to name the object, the operation, and the assumption before using the compact notation.



\paragraph{Step-by-step derivation.}

For Curved Spacetime and the Standard Model: Boundary conditions, normalization, and signs: identifying physical fields, the derivation is written from the smallest usable model upward.

List the fields with their representations under \(SU(3)_c\times SU(2)_L\times U(1)_Y\).  A term in the Lagrangian is allowed only if all gauge indices can be contracted to a singlet and the hypercharges sum to zero.  For example, the charged-lepton Yukawa term has the schematic form
\[
  -y_e\bar L\Phi e_R+\hbox{h.c.}
\]
The \(SU(2)\) indices of \(L\) and \(\Phi\) are contracted, and the hypercharges are chosen so the product is neutral.  When \(\Phi\) is expanded around its vacuum value, the same term becomes a fermion mass term.  This single check uses representations, gauge invariance, spontaneous symmetry breaking, and ordinary dimensional analysis.  In Curved Spacetime and the Standard Model, section 6, the useful habit is to name the object, the operation, and the assumption before using the compact notation.

The formulas to keep in view are

\[
  G_{\mathrm{SM}}=SU(3)_c\times SU(2)_L\times U(1)_Y
\]
\[
  D_\mu=\p_\mu+\ii g_sG_\mu^aT^a+\ii gW_\mu^i\tau^i+\ii g'YB_\mu
\]
\[
  \Lag\supset -y_e\bar L\Phi e_R+\hbox{h.c.}
\]

In this setting the calculation for Curved Spacetime and the Standard Model: Boundary conditions, normalization, and signs: identifying physical fields is complete when the finite model, the limiting step, and the normalization or boundary condition can all be named without looking back at the prose.




\begin{example}[A worked calculation for Curved Spacetime and the Standard Model: Boundary conditions, normalization, and signs: identifying physical fields]
Check a charged-lepton Yukawa term.  The left-handed lepton \(L\) is an \(SU(2)\) doublet, the Higgs field \(\Phi\) is a doublet, and \(e_R\) is a singlet.  Contract the two doublet indices and choose hypercharges so the sum is zero.  After \(\Phi\) gets a vacuum value, the term becomes \(m_e\bar e e\) with \(m_e=y_ev/\sqrt2\).  In Curved Spacetime and the Standard Model, section 6, the useful habit is to name the object, the operation, and the assumption before using the compact notation.
\end{example}



\begin{exercise}
State the Standard Model gauge group.  Use the notation of Curved Spacetime and the Standard Model: Boundary conditions, normalization, and signs: identifying physical fields.
\begin{answer}
\(SU(3)_c\times SU(2)_L\times U(1)_Y\).
\end{answer}
\end{exercise}

\begin{exercise}
Why must hypercharges in an interaction term sum to zero?  Use the notation of Curved Spacetime and the Standard Model: Boundary conditions, normalization, and signs: identifying physical fields.
\begin{answer}
Otherwise the term picks up a local \(U(1)_Y\) phase and is not gauge invariant.
\end{answer}
\end{exercise}



\section{How the idea enters quantum field theory: placing QFT on a background metric}

The work of this section is placing QFT on a background metric.  In the chapter heading the vocabulary is curved backgrounds, Standard Model, electroweak theory, but the first objects on the page are more prosaic: gauge groups, matter multiplets, Yukawa couplings, electroweak breaking, curved backgrounds, and effective descriptions.  The formal notation will be useful only after it has been tied to a limit, a symmetry, a normalization condition, or an integration-by-parts step.  In Curved Spacetime and the Standard Model, section 7, the useful habit is to name the object, the operation, and the assumption before using the compact notation.

The finite model is a table listing fields, representations, and charges, checked term by term for gauge invariance.  In that model no mystery is available: every derivative is an ordinary derivative with respect to one listed variable, every inner product is a sum, and every boundary assumption can be written at the endpoints.  This is why the finite model is not a toy in the dismissive sense.  It is the bookkeeping device that prevents continuum notation from becoming a fog of indices.  For Curved Spacetime and the Standard Model, section 7, that bookkeeping is deliberately modest, but it is what later keeps signs, factors of \(2\pi\), and normalization constants from turning into guesses.

The central idea for this chapter is the Standard Model Lagrangian is the collection of lowest-dimension local terms allowed by its gauge symmetries and field content.  To test that idea, strip the formula down until only the operation remains.  A sign change after raising or lowering an index means the metric has entered.  A derivative moving from one factor to another means a boundary term has been handled.  A normalization constant means that some unit object, such as a delta function, basis vector, or one-particle state, has been fixed.  Read in the setting of Curved Spacetime and the Standard Model, section 7, the line is a check on the calculation rather than an invitation to memorize another rule.

For how the idea enters quantum field theory: placing qft on a background metric, the calculation should also pass a dimensional check.  The left and right sides must carry the same units; more subtly, they must transform in the same way under the symmetries already introduced.  This is not a proof, but it is a quick way to catch wrong factors of \(2\pi\), signs in the metric, missing powers of the lattice spacing, or an omitted charge.  The same step will look more abstract later; in Curved Spacetime and the Standard Model, section 7 it is kept close to the finite calculation so the limiting move remains visible.

The bridge to quantum field theory is direct.  The final application uses every earlier tool: spinors, gauge fields, symmetry breaking, renormalization, and effective reasoning.  Later notation will carry more indices and fewer verbal reminders, so the safe habit is to keep asking which operation is being performed and which quantity is being held fixed.  At this point in Curved Spacetime and the Standard Model, section 7, it is worth asking what would fail if the boundary condition, normalization, or symmetry assumption were changed.

One warning is worth making explicit here.  A compact Lagrangian hides many definitions; each symbol carries representation, charge, and normalization data.  This warning points to the delicate step of the derivation, not to a decorative aside.  In Curved Spacetime and the Standard Model, section 7, the calculation is meant to leave a trail from an ordinary derivative, sum, matrix product, or Gaussian integral to the field-theory formula.

The section closes by keeping placing QFT on a background metric connected to Curved Spacetime and the Standard Model.  The same general operation will be used with different objects in neighboring chapters, so the present calculation is both a result and a rehearsal.  In Curved Spacetime and the Standard Model, section 7, the useful habit is to name the object, the operation, and the assumption before using the compact notation.



\paragraph{Step-by-step derivation.}

For Curved Spacetime and the Standard Model: How the idea enters quantum field theory: placing QFT on a background metric, the derivation is written from the smallest usable model upward.

List the fields with their representations under \(SU(3)_c\times SU(2)_L\times U(1)_Y\).  A term in the Lagrangian is allowed only if all gauge indices can be contracted to a singlet and the hypercharges sum to zero.  For example, the charged-lepton Yukawa term has the schematic form
\[
  -y_e\bar L\Phi e_R+\hbox{h.c.}
\]
The \(SU(2)\) indices of \(L\) and \(\Phi\) are contracted, and the hypercharges are chosen so the product is neutral.  When \(\Phi\) is expanded around its vacuum value, the same term becomes a fermion mass term.  This single check uses representations, gauge invariance, spontaneous symmetry breaking, and ordinary dimensional analysis.  In Curved Spacetime and the Standard Model, section 7, the useful habit is to name the object, the operation, and the assumption before using the compact notation.

The formulas to keep in view are

\[
  G_{\mathrm{SM}}=SU(3)_c\times SU(2)_L\times U(1)_Y
\]
\[
  D_\mu=\p_\mu+\ii g_sG_\mu^aT^a+\ii gW_\mu^i\tau^i+\ii g'YB_\mu
\]
\[
  \Lag\supset -y_e\bar L\Phi e_R+\hbox{h.c.}
\]

In this setting the calculation for Curved Spacetime and the Standard Model: How the idea enters quantum field theory: placing QFT on a background metric is complete when the finite model, the limiting step, and the normalization or boundary condition can all be named without looking back at the prose.




\begin{example}[A worked calculation for Curved Spacetime and the Standard Model: How the idea enters quantum field theory: placing QFT on a background metric]
Check a charged-lepton Yukawa term.  The left-handed lepton \(L\) is an \(SU(2)\) doublet, the Higgs field \(\Phi\) is a doublet, and \(e_R\) is a singlet.  Contract the two doublet indices and choose hypercharges so the sum is zero.  After \(\Phi\) gets a vacuum value, the term becomes \(m_e\bar e e\) with \(m_e=y_ev/\sqrt2\).  In Curved Spacetime and the Standard Model, section 7, the useful habit is to name the object, the operation, and the assumption before using the compact notation.
\end{example}



\begin{exercise}
State the Standard Model gauge group.  Use the notation of Curved Spacetime and the Standard Model: How the idea enters quantum field theory: placing QFT on a background metric.
\begin{answer}
\(SU(3)_c\times SU(2)_L\times U(1)_Y\).
\end{answer}
\end{exercise}

\begin{exercise}
Why must hypercharges in an interaction term sum to zero?  Use the notation of Curved Spacetime and the Standard Model: How the idea enters quantum field theory: placing QFT on a background metric.
\begin{answer}
Otherwise the term picks up a local \(U(1)_Y\) phase and is not gauge invariant.
\end{answer}
\end{exercise}



\section{Review problems and extensions: summarizing the full course}

The work of this section is summarizing the full course.  In the chapter heading the vocabulary is curved backgrounds, Standard Model, electroweak theory, but the first objects on the page are more prosaic: gauge groups, matter multiplets, Yukawa couplings, electroweak breaking, curved backgrounds, and effective descriptions.  The formal notation will be useful only after it has been tied to a limit, a symmetry, a normalization condition, or an integration-by-parts step.  In Curved Spacetime and the Standard Model, section 8, the useful habit is to name the object, the operation, and the assumption before using the compact notation.

The finite model is a table listing fields, representations, and charges, checked term by term for gauge invariance.  In that model no mystery is available: every derivative is an ordinary derivative with respect to one listed variable, every inner product is a sum, and every boundary assumption can be written at the endpoints.  This is why the finite model is not a toy in the dismissive sense.  It is the bookkeeping device that prevents continuum notation from becoming a fog of indices.  For Curved Spacetime and the Standard Model, section 8, that bookkeeping is deliberately modest, but it is what later keeps signs, factors of \(2\pi\), and normalization constants from turning into guesses.

The central idea for this chapter is the Standard Model Lagrangian is the collection of lowest-dimension local terms allowed by its gauge symmetries and field content.  To test that idea, strip the formula down until only the operation remains.  A sign change after raising or lowering an index means the metric has entered.  A derivative moving from one factor to another means a boundary term has been handled.  A normalization constant means that some unit object, such as a delta function, basis vector, or one-particle state, has been fixed.  Read in the setting of Curved Spacetime and the Standard Model, section 8, the line is a check on the calculation rather than an invitation to memorize another rule.

For review problems and extensions: summarizing the full course, the calculation should also pass a dimensional check.  The left and right sides must carry the same units; more subtly, they must transform in the same way under the symmetries already introduced.  This is not a proof, but it is a quick way to catch wrong factors of \(2\pi\), signs in the metric, missing powers of the lattice spacing, or an omitted charge.  The same step will look more abstract later; in Curved Spacetime and the Standard Model, section 8 it is kept close to the finite calculation so the limiting move remains visible.

The bridge to quantum field theory is direct.  The final application uses every earlier tool: spinors, gauge fields, symmetry breaking, renormalization, and effective reasoning.  Later notation will carry more indices and fewer verbal reminders, so the safe habit is to keep asking which operation is being performed and which quantity is being held fixed.  At this point in Curved Spacetime and the Standard Model, section 8, it is worth asking what would fail if the boundary condition, normalization, or symmetry assumption were changed.

One warning is worth making explicit here.  A compact Lagrangian hides many definitions; each symbol carries representation, charge, and normalization data.  This warning points to the delicate step of the derivation, not to a decorative aside.  In Curved Spacetime and the Standard Model, section 8, the calculation is meant to leave a trail from an ordinary derivative, sum, matrix product, or Gaussian integral to the field-theory formula.

The section closes by keeping summarizing the full course connected to Curved Spacetime and the Standard Model.  The same general operation will be used with different objects in neighboring chapters, so the present calculation is both a result and a rehearsal.  In Curved Spacetime and the Standard Model, section 8, the useful habit is to name the object, the operation, and the assumption before using the compact notation.



\paragraph{Step-by-step derivation.}

For Curved Spacetime and the Standard Model: Review problems and extensions: summarizing the full course, the derivation is written from the smallest usable model upward.

List the fields with their representations under \(SU(3)_c\times SU(2)_L\times U(1)_Y\).  A term in the Lagrangian is allowed only if all gauge indices can be contracted to a singlet and the hypercharges sum to zero.  For example, the charged-lepton Yukawa term has the schematic form
\[
  -y_e\bar L\Phi e_R+\hbox{h.c.}
\]
The \(SU(2)\) indices of \(L\) and \(\Phi\) are contracted, and the hypercharges are chosen so the product is neutral.  When \(\Phi\) is expanded around its vacuum value, the same term becomes a fermion mass term.  This single check uses representations, gauge invariance, spontaneous symmetry breaking, and ordinary dimensional analysis.  In Curved Spacetime and the Standard Model, section 8, the useful habit is to name the object, the operation, and the assumption before using the compact notation.

The formulas to keep in view are

\[
  G_{\mathrm{SM}}=SU(3)_c\times SU(2)_L\times U(1)_Y
\]
\[
  D_\mu=\p_\mu+\ii g_sG_\mu^aT^a+\ii gW_\mu^i\tau^i+\ii g'YB_\mu
\]
\[
  \Lag\supset -y_e\bar L\Phi e_R+\hbox{h.c.}
\]

In this setting the calculation for Curved Spacetime and the Standard Model: Review problems and extensions: summarizing the full course is complete when the finite model, the limiting step, and the normalization or boundary condition can all be named without looking back at the prose.




\begin{example}[A worked calculation for Curved Spacetime and the Standard Model: Review problems and extensions: summarizing the full course]
Check a charged-lepton Yukawa term.  The left-handed lepton \(L\) is an \(SU(2)\) doublet, the Higgs field \(\Phi\) is a doublet, and \(e_R\) is a singlet.  Contract the two doublet indices and choose hypercharges so the sum is zero.  After \(\Phi\) gets a vacuum value, the term becomes \(m_e\bar e e\) with \(m_e=y_ev/\sqrt2\).  In Curved Spacetime and the Standard Model, section 8, the useful habit is to name the object, the operation, and the assumption before using the compact notation.
\end{example}



\begin{exercise}
State the Standard Model gauge group.  Use the notation of Curved Spacetime and the Standard Model: Review problems and extensions: summarizing the full course.
\begin{answer}
\(SU(3)_c\times SU(2)_L\times U(1)_Y\).
\end{answer}
\end{exercise}

\begin{exercise}
Why must hypercharges in an interaction term sum to zero?  Use the notation of Curved Spacetime and the Standard Model: Review problems and extensions: summarizing the full course.
\begin{answer}
Otherwise the term picks up a local \(U(1)_Y\) phase and is not gauge invariant.
\end{answer}
\end{exercise}



\section{Cumulative worked derivation: from reading the gauge group to placing QFT on a background metric}

This final worked section braids together reading the gauge group, checking Yukawa terms, and placing QFT on a background metric for Curved Spacetime and the Standard Model.  The vocabulary of the chapter was curved backgrounds, Standard Model, electroweak theory; here those words are used in one continuous calculation rather than in isolated fragments.

Read the Standard Model as a worked exercise in allowed local terms.  The gauge group fixes the covariant derivatives.  The matter representations decide which bilinears can be formed.  The Higgs representation allows Yukawa terms and breaks \(SU(2)_L\times U(1)_Y\) to electromagnetism.  Renormalization explains why the same Lagrangian can be used at different scales, while effective operators record possible short-distance physics not resolved by the theory.  In Curved Spacetime and the Standard Model, cumulative derivation, the useful habit is to name the object, the operation, and the assumption before using the compact notation.

For reference, the compact formulas are

\[
  G_{\mathrm{SM}}=SU(3)_c\times SU(2)_L\times U(1)_Y
\]
\[
  D_\mu=\p_\mu+\ii g_sG_\mu^aT^a+\ii gW_\mu^i\tau^i+\ii g'YB_\mu
\]
\[
  \Lag\supset -y_e\bar L\Phi e_R+\hbox{h.c.}
\]

\begin{exercise}
Reproduce the cumulative derivation for Curved Spacetime and the Standard Model without looking at the prose.  Mark the step where a finite calculation becomes a continuum, operator, or field-theoretic statement.
\begin{answer}
The reconstruction should name the starting object, carry out the differentiating, varying, transforming, or commuting step explicitly, and state the normalization or boundary condition that makes the final formula meaningful.
\end{answer}
\end{exercise}



\section{Chapter synthesis}

This chapter has treated curved spacetime and the standard model as a chain of calculations.  The safest summary is not a list of final formulas, but a map from assumptions to equations: finite model, limiting step, boundary condition, normalization, and field-theory use.  If one of those links is missing, the formula should be rederived before it is used in a later chapter.

\begin{exercise}
Write a one-page reconstruction of curved spacetime and the standard model.  Include one equation from the finite model, one equation from the continuum or operator form, and one sentence explaining how the two are connected.
\begin{answer}
A strong reconstruction names the basic objects, identifies the central derivative, variation, commutator, or symmetry operation, and states the assumption that allows the finite calculation to become the field-theory formula.
\end{answer}
\end{exercise}
